%
%      Brno University of Technology
%      Faculty of Information Technology
%
%      MSc Thesis	2009/2010
%
% 		 Design of S-Boxes Using Genetic Algorithms
%
%
%      Bedrich Hovorka
%
%V této práci příliš žluťoučký kůň úpěl ďábelské ódy gg
Rostoucí výkon počítačů a nárůst jejich použití širokou masou uživatelů v celosvětové po\-čí\-ta\-čo\-vé síti
si žádají větší pozornost věnovanou bezpečnosti těchto uživatelů.
Neboť součástí této sítě
nejsou jen samí lidé a organizace s dobrými úmysly, ale i tací, kteří se snaží
porušit právo na soukromí odposlechnutím důvěrných
informací jiné osoby a využít je ve svůj prospěch či tu osobu zkompromitovat jiným způsobem, například
se za ni vydávat. Proto je důležité těmto záškodníkům tyto činnosti prakticky znemožnit převodem důvěrné informace
do podoby, kterou budou schopny přečíst jen vybrané osoby a nebude ji schopna přečíst žádná nežádoucí osoba
-- důvěrnost.
Dále musí být součástí tohoto převodu ověření identity původce přenášené informace -- autentizace.
Algoritmy, které zajišťují důvěrnost a autentizaci přenášených dat, zkoumá kryptografie.
Těmto algoritmům se též říká šifry. Díky čím dál tím většímu používání šifer a objevování novějších
druhů kryptoanalýzy jejich životnost klesá.
Poučení z historie šifer říká, že šifra byla neprolomitelná do té doby, než
někdo objevil její slabinu. Proto musela být vzápětí vždy vytvořena šifra nová.
V dnešní době probíhá neustálý výzkum a vývoj nových šifer, které by měly být odolnější
vůči stále přibývajícím novějším druhům útoků.

Tato práce se zabývá částí šifrovacího algoritmu zvanou substituční box
(zkráceně s-box) a jeho vývojem. K jejich konstrukci se využívají jak čistě
matematické přístupy s využitím teorie Galoisových polí, tak i různé stochastické způsoby návrhu, ke
kterým se řadí například evoluce, která je použita v této práci.
Evoluční výpočetní techniky jsou, jak již název napovídá, inspirovány evoluční teorií z biologie.
Cílem práce je použít některé vlastnosti substitučních boxů
jako objektivní funkce k jejich evolučnímu vývoji.
Pře\-sto\-že již bylo mnoho vlastností studováno, relativně málo práce bylo uskutečněno
při hledání s-boxů především v multikriteriálním genetickém prohledávání.

Nejprve však bylo nutné navrhnout a naimplementovat kolekci algoritmů,
která tyto vlastnosti umožní experimentálně zjistit.
Tato implementace musí především umožnit rychle zadávat různé evoluční experimenty a pokud možno
provést je co nejrychleji.

Tato práce vznikla jako samostatné dílo, které nenavazuje na žádný školní projekt.
